% ============================================================
% ANEXO B · SALIDAS Y RESULTADOS COMPLEMENTARIOS
% Tablas generadas automáticamente a partir de los datos en Code/results/.
% ============================================================

\section{Salidas y resultados complementarios}\label{anx:salidas}

En este anexo se presentan las tablas cuantitativas complementarias que respaldan las
etapas de preprocesamiento, selección de variables, validación cruzada, coste computacional
y verificación de consistencia interna del sistema. Todas las tablas han sido generadas
directamente a partir de los registros estructurados en archivos CSV del repositorio
mediante el script \texttt{tablas\_anexos.py} de \texttt{Code/tools/}, garantizando trazabilidad completa.

% ------------------------------------------------------------
\subsection{Preprocesamiento y auditoría de máscaras (Módulo 2)}

La Tabla~\ref{tab:anexo_m2_resumen} sintetiza el control de calidad aplicado sobre las 3\,568
mamografías, documentando los casos con desajuste dimensional corregidos y la concordancia con
los recortes oficiales.

% Generado por Code/tools/tablas_anexos.py desde 2_preprocesamiento_resumen.csv. No editar a mano.
\begin{table}[H]
\centering
\small
\caption{Resumen del preprocesamiento, alineación de máscaras y control de calidad (Módulo 2).}
\label{tab:anexo_m2_resumen}
\begin{adjustbox}{max width=\linewidth}
\begin{tabular}{l r r r r r c r}
\toprule
Subconjunto & Lesiones & Desalineadas & Máscaras & Recortes & Recortes oficiales & Dimensión mediana & Recortes \\
 & analizadas & (escal./offset) & depuradas & coincidentes & con discrepancia/sin & de lesión (px) & en borde \\
\midrule
Calcificaciones & 1872 & 2 & 0 & 1663 & 209 & 273$\times$289 & 120 \\
Masas & 1696 & 78 & 1464 & 1543 & 75 & 302$\times$304 & 67 \\
\bottomrule
\end{tabular}
\end{adjustbox}
\par\smallskip
\parbox{0.95\linewidth}{\footnotesize Auditoría y control de calidad de M2. Desalineadas: 78 por escala ($0.870$) en masas y 2 por traslación en calcificaciones. Depuradas: reducción a la componente conexa mayor eliminando artefactos menores aislados. Recortes coincidentes: concordancia exacta al píxel con los recortes oficiales del CBIS-DDSM.}
\end{table}


% ------------------------------------------------------------
\subsection{Selección y redundancia de características (Módulo 4)}

La Tabla~\ref{tab:anexo_m4_seleccion} reporta las 12 características radiómicas seleccionadas
en cada subconjunto clínico, acompañadas de sus estadísticos univariados de significación y tamaño del efecto.

% Generado por Code/tools/tablas_anexos.py desde 4_seleccion.csv y 3_tamano_efecto.csv. No editar a mano.
\begin{table}[H]
\centering
\small
\caption{Conjunto completo de características radiómicas seleccionadas por subconjunto (Módulo 4).}
\label{tab:anexo_m4_seleccion}
\begin{adjustbox}{max width=\linewidth}
\begin{tabular}{r l l r r r l l r r r}
\toprule
 & \multicolumn{5}{c}{Masas} & \multicolumn{5}{c}{Calcificaciones} \\ \cmidrule(lr){2-6}\cmidrule(lr){7-11}
\# & Característica & Familia & $F$ & $p$ & $|d|$ & Característica & Familia & $F$ & $p$ & $|d|$ \\
\midrule
1 & firstorder\_Maximum & firstorder & 83.1 & $<$0.001 & 0.502 & shape2D\_MinorAxisLength & shape2D & 235.4 & $<$0.001 & 0.817 \\
2 & firstorder\_90Percentile & firstorder & 82.6 & $<$0.001 & 0.501 & glrlm\_RunEntropy & glrlm & 200.9 & $<$0.001 & 0.755 \\
3 & firstorder\_Energy & firstorder & 74.7 & $<$0.001 & 0.476 & shape2D\_Perimeter & shape2D & 177.0 & $<$0.001 & 0.709 \\
4 & shape2D\_MinorAxisLength & shape2D & 66.7 & $<$0.001 & 0.450 & firstorder\_Skewness & firstorder & 174.6 & $<$0.001 & 0.704 \\
5 & shape2D\_MaximumDiameter & shape2D & 62.3 & $<$0.001 & 0.435 & glrlm\_ShortRunLowGrayLevelEmphasis & glrlm & 166.2 & $<$0.001 & 0.687 \\
6 & shape2D\_Perimeter & shape2D & 53.1 & $<$0.001 & 0.402 & glrlm\_ShortRunEmphasis & glrlm & 161.5 & $<$0.001 & 0.677 \\
7 & shape2D\_MeshSurface & shape2D & 52.8 & $<$0.001 & 0.400 & glcm\_JointAverage & glcm & 150.9 & $<$0.001 & 0.654 \\
8 & shape2D\_PerimeterSurfaceRatio & shape2D & 51.3 & $<$0.001 & 0.395 & glcm\_InverseVariance & glcm & 142.8 & $<$0.001 & 0.636 \\
9 & glrlm\_GrayLevelNonUniformity & glrlm & 37.8 & $<$0.001 & 0.339 & glcm\_DifferenceEntropy & glcm & 135.6 & $<$0.001 & 0.620 \\
10 & firstorder\_10Percentile & firstorder & 31.7 & $<$0.001 & 0.311 & firstorder\_Energy & firstorder & 135.1 & $<$0.001 & 0.619 \\
11 & glcm\_Correlation & glcm & 31.0 & $<$0.001 & 0.307 & glcm\_ClusterShade & glcm & 133.9 & $<$0.001 & 0.616 \\
12 & glrlm\_RunLengthNonUniformity & glrlm & 28.8 & $<$0.001 & 0.296 & shape2D\_PixelSurface & shape2D & 129.1 & $<$0.001 & 0.605 \\
\bottomrule
\end{tabular}
\end{adjustbox}
\par\smallskip
\parbox{0.95\linewidth}{\footnotesize Selección de las 12 características en cada subconjunto mediante ANOVA univariado ($F$-test) sobre el conjunto de entrenamiento con restricción de redundancia lineal $|r| \le 0.95$. $p$ reporta el nivel de significación estadística; $|d|$ es el tamaño del efecto (d de Cohen sobre intensidades crudas).}
\end{table}


La Tabla~\ref{tab:anexo_m4_redundancia} documenta las características candidatas que, a pesar de su
relevancia individual, fueron descartadas del conjunto final por exceder el límite de redundancia lineal.

% Generado por Code/tools/tablas_anexos.py desde 4_descartadas_redundancia.csv. No editar a mano.
\begin{table}[H]
\centering
\small
\caption{Características radiómicas descartadas por restricción de redundancia (Módulo 4).}
\label{tab:anexo_m4_redundancia}
\begin{adjustbox}{max width=\linewidth}
\begin{tabular}{l l r l}
\toprule
Característica candidata & Redundante con & $|r|$ & Estado \\
\midrule
\multicolumn{4}{l}{\textit{Masas}} \\
\quad firstorder\_TotalEnergy & firstorder\_Energy & 1.000 & Descartada ($|r| > 0.95$) \\
\quad firstorder\_Median & firstorder\_90Percentile & 0.959 & Descartada ($|r| > 0.95$) \\
\quad firstorder\_RootMeanSquared & firstorder\_90Percentile & 0.966 & Descartada ($|r| > 0.95$) \\
\quad firstorder\_Mean & firstorder\_90Percentile & 0.955 & Descartada ($|r| > 0.95$) \\
\quad shape2D\_MajorAxisLength & shape2D\_MaximumDiameter & 0.993 & Descartada ($|r| > 0.95$) \\
\quad shape2D\_PixelSurface & shape2D\_MeshSurface & 1.000 & Descartada ($|r| > 0.95$) \\
\quad glcm\_MCC & glcm\_Correlation & 0.986 & Descartada ($|r| > 0.95$) \\
\multicolumn{4}{l}{\textit{Calcificaciones}} \\
\quad glcm\_Imc1 & glrlm\_RunEntropy & 0.961 & Descartada ($|r| > 0.95$) \\
\quad glcm\_SumAverage & glcm\_JointAverage & 1.000 & Descartada ($|r| > 0.95$) \\
\quad glrlm\_RunPercentage & glrlm\_ShortRunEmphasis & 0.975 & Descartada ($|r| > 0.95$) \\
\quad glrlm\_RunLengthNonUniformityNormalized & glrlm\_ShortRunEmphasis & 0.980 & Descartada ($|r| > 0.95$) \\
\quad shape2D\_MaximumDiameter & shape2D\_Perimeter & 0.980 & Descartada ($|r| > 0.95$) \\
\quad firstorder\_TotalEnergy & firstorder\_Energy & 1.000 & Descartada ($|r| > 0.95$) \\
\quad glrlm\_HighGrayLevelRunEmphasis & glcm\_JointAverage & 0.972 & Descartada ($|r| > 0.95$) \\
\quad glcm\_Autocorrelation & glcm\_JointAverage & 0.981 & Descartada ($|r| > 0.95$) \\
\bottomrule
\end{tabular}
\end{adjustbox}
\par\smallskip
\parbox{0.95\linewidth}{\footnotesize Características candidatas con alto valor de $F$ que fueron omitidas de la selección final para evitar colinealidad extrema ($|r| > 0.95$). En masas se descartaron 7 variables y en calcificaciones 8.}
\end{table}


% ------------------------------------------------------------
\subsection{Estratificación de pliegues en validación cruzada (Módulo 4)}

La Tabla~\ref{tab:anexo_m4_folds} detalla el reparto de lesiones, el número de pacientes únicos y
la proporción de malignidad en cada uno de los cinco pliegues de validación cruzada generados.

% Generado por Code/tools/tablas_anexos.py desde 4_folds.csv. No editar a mano.
\begin{table}[H]
\centering
\small
\caption{Balance de lesiones, pacientes y proporción de malignidad por fold (Módulo 4).}
\label{tab:anexo_m4_folds}
\begin{adjustbox}{max width=\linewidth}
\begin{tabular}{c r r r r r r}
\toprule
 & \multicolumn{3}{c}{Masas (semilla elegida: 125)} & \multicolumn{3}{c}{Calcificaciones (semilla elegida: 168)} \\ \cmidrule(lr){2-4}\cmidrule(lr){5-7}
Fold & Lesiones & Pacientes & Malignas (\%) & Lesiones & Pacientes & Malignas (\%) \\
\midrule
0 & 255 & 139 & 48.6\% & 315 & 120 & 33.3\% \\
1 & 258 & 134 & 47.3\% & 271 & 119 & 35.4\% \\
2 & 268 & 138 & 48.9\% & 322 & 121 & 36.6\% \\
3 & 257 & 140 & 49.4\% & 314 & 121 & 35.7\% \\
4 & 280 & 140 & 47.5\% & 324 & 121 & 34.9\% \\
\midrule
Total / media & 1318 & 691 & 48.3\% & 1546 & 602 & 35.2\% \\
\bottomrule
\end{tabular}
\end{adjustbox}
\par\smallskip
\parbox{0.95\linewidth}{\footnotesize Partición de validación cruzada 5-fold estratificada por paciente (\texttt{StratifiedGroupKFold}). Garantiza cero fuga de pacientes entre folds y balance de clases óptimo entre candidatos de semillas.}
\end{table}


% ------------------------------------------------------------
\subsection{Costes computacionales de simulación cuántica (Módulos 5, 8 y X8)}

La Tabla~\ref{tab:anexo_m5_tiempos} presenta la comparativa de tiempos de ejecución empíricos entre
el cálculo matricial por vector de estado y la simulación individual por circuito.

% Generado por Code/tools/tablas_anexos.py desde 5b_tiempos.csv. No editar a mano.
\begin{table}[H]
\centering
\small
\caption{Costes computacionales de evaluación de kernels cuánticos (Módulo 5).}
\label{tab:anexo_m5_tiempos}
\begin{adjustbox}{max width=\linewidth}
\begin{tabular}{l l c r}
\toprule
Subconjunto & Kernel cuántico evaluado & Dimensión & Tiempo total (s) \\
\midrule
Masas & \texttt{FidelityStatevectorKernel\_C4} & 200x200 & 1.468 \\
Masas & \texttt{FidelityStatevectorKernel\_C5} & 200x200 & 1.590 \\
Calcificaciones & \texttt{FidelityStatevectorKernel\_C4} & 200x200 & 1.569 \\
Calcificaciones & \texttt{FidelityStatevectorKernel\_C5} & 200x200 & 1.574 \\
Masas & \texttt{FidelityQuantumKernel\_C4} & 10x10 & 4.970 \\
Masas & \texttt{FidelityQuantumKernel\_C5} & 10x10 & 4.985 \\
Calcificaciones & \texttt{FidelityQuantumKernel\_C4} & 10x10 & 4.948 \\
Calcificaciones & \texttt{FidelityQuantumKernel\_C5} & 10x10 & 5.041 \\
\bottomrule
\end{tabular}
\end{adjustbox}
\par\smallskip
\parbox{0.95\linewidth}{\footnotesize Tiempos medidos en simulación exacta. \texttt{FidelityStatevectorKernel} calcula la matriz $200\times 200$ en $\sim$1.5\,s aprovechando vectores de estado precomputados. \texttt{FidelityQuantumKernel} evalúa el circuito cuántico directo requiriendo $\sim$5.0\,s para una matriz reducida de solo $10\times 10$.}
\end{table}


La Tabla~\ref{tab:anexo_m8_oe6_qubits} reporta el escalamiento de los requerimientos de memoria y tiempo
de simulación en función del número de qubits del sistema, contrastando mediciones directas y extrapolaciones.

% Generado por Code/tools/tablas_anexos.py desde 8_oe6_qubits.csv. No editar a mano.
\begin{table}[H]
\centering
\small
\caption{Escalamiento computacional con el número de qubits en simulación por vector de estado (OE-6).}
\label{tab:anexo_m8_oe6_qubits}
\begin{adjustbox}{max width=\linewidth}
\begin{tabular}{c c r r r r}
\toprule
$k$ (qubits) & Tipo de dato & Memoria estado & Tiempo/estado (s) & Tiempo embedding (h) & Entrenamiento (h/mapa) \\
\midrule
8 & Medido & 0.00 MiB & 0.0012 & 0.002 & 1.6 \\
10 & Medido & 0.02 MiB & 0.0019 & 0.004 & 3.0 \\
12 & Medido & 0.06 MiB & 0.0029 & 0.006 & 5.7 \\
14 & Medido & 0.25 MiB & 0.0065 & 0.013 & 14.8 \\
16 & Medido & 1.00 MiB & 0.0349 & 0.069 & 90.3 \\
18 & Medido & 4.00 MiB & 0.1071 & 0.212 & 311 \\
20 & Medido & 16.00 MiB & 0.3924 & 0.778 & 1264 \\
22 & Medido & 64.00 MiB & 1.2 & 2.472 & 4415 \\
24 & Medido & 256.00 MiB & 4.1 & 8.194 & 15950 \\
26 & Medido & 1.00 GiB & 13.7 & 27.3 & 57421 \\
28 & Medido & 4.00 GiB & 74.8 & 148.2 & 336114 \\
30 & Extrapolado & 16.00 GiB & 221.0 & 438.0 & 1063514 \\
32 & Extrapolado & 64.00 GiB & 802.7 & 1591.1 & 4118952 \\
\bottomrule
\end{tabular}
\end{adjustbox}
\par\smallskip
\parbox{0.95\linewidth}{\footnotesize Escalamiento del coste en simulación cuántica por vector de estado en función del número de qubits ($k$). Para $k \le 28$ las cifras corresponden a ejecuciones empíricas en el hardware de referencia (Apple M3); para $k \ge 30$ corresponden a la extrapolación exponencial teórica.}
\end{table}


La Tabla~\ref{tab:anexo_x8_gradiente} expone la evaluación comparativa del coste por muestra para el
cálculo de gradientes cuánticos mediante la regla analítica de desplazamiento de parámetros frente a diferenciación automática.

% Generado por Code/tools/tablas_anexos.py desde X8_costo_gradiente.csv. No editar a mano.
\begin{table}[H]
\centering
\small
\caption{Coste computacional del cálculo de gradientes variacionales (Experimento X8).}
\label{tab:anexo_x8_gradiente}
\begin{adjustbox}{max width=\linewidth}
\begin{tabular}{l l r r}
\toprule
Método de cálculo & Naturaleza del gradiente & Tiempo/lesión (s) & Coste estimado (50 épocas, h) \\
\midrule
parameter shift (EstimatorQNN, exact) & Jacobian 12 $\times$ 24 (also needed for a scalar loss) & 5.1589 & 205.21 \\
automatic differentiation (state vector) & Jacobian 12 $\times$ 24 & 0.2369 & 9.42 \\
automatic differentiation (state vector) & scalar loss, one backward pass & 0.0014 & 0.06 \\
\bottomrule
\end{tabular}
\end{adjustbox}
\par\smallskip
\parbox{0.95\linewidth}{\footnotesize Comparación de costes entre diferenciación exacta por \textit{parameter-shift rule} y diferenciación automática directa sobre vector de estado en simulación (Experimento X8).}
\end{table}


% ------------------------------------------------------------
\subsection{Verificaciones automáticas y control metodológico (Módulos 5 a 8 y X5)}

La Tabla~\ref{tab:anexo_verificaciones} resume el catálogo de verificaciones automáticas integradas
en los módulos de transformación, separabilidad y clasificación supervisada.

% Generado por Code/tools/tablas_anexos.py desde 5a_, 5b_, 6_, 7_ y 8_verificaciones.csv. No editar a mano.
\begin{table}[H]
\centering
\small
\caption{Resumen de la suite de verificaciones automáticas de consistencia y correctitud.}
\label{tab:anexo_verificaciones}
\begin{adjustbox}{max width=\linewidth}
\begin{tabular}{l l c c p{6.5cm}}
\toprule
Módulo & Archivo de resultados & Pruebas & Superadas & Ámbito de verificación \\
\midrule
Módulo 5a & \texttt{5a\_verificaciones.csv} & 18 & 18 (100\%) & Recomputación exacta, dispersión positiva y validación con SparsePauliOp \\
Módulo 5b & \texttt{5b\_verificaciones.csv} & 38 & 38 (100\%) & Simetría de matrices, diagonal unitaria y semidefinición positiva ($\lambda_{\min} \ge 0$) \\
Módulo 6 & \texttt{6\_verificaciones.csv} & 20 & 20 (100\%) & Normalización de traza, identidades teóricas de $g$, KTA y control nulo C1=C2 \\
Módulo 7 & \texttt{7\_verificaciones.csv} & 16 & 16 (100\%) & Cero fuga de pacientes, convergencia por parada temprana y estabilidad de probabilidades \\
Módulo 8 & \texttt{8\_verificaciones.csv} & 22 & 22 (100\%) & Reproducción de AUC en test y CV, nula exacta de Spearman y muestreo de $\langle Z_i \rangle$ \\
\midrule
Total suite & 5 archivos CSV & 114 & 114 (100\%) & Suite completa de aserciones automáticas \\
\bottomrule
\end{tabular}
\end{adjustbox}
\par\smallskip
\parbox{0.95\linewidth}{\footnotesize Todas las pruebas automáticas del pipeline ejecutan aserciones matemáticas estrictas sobre las estructuras de datos generadas. El 100\,\% de las verificaciones resultó exitoso, sin discrepancias numéricas tolerables más allá del error de redondeo de punto flotante.}
\end{table}


La Tabla~\ref{tab:anexo_x5_auditoria} reúne los resultados de la auditoría independiente de consistencia
metodológica, verificando la ausencia de fugas de datos y la integridad de las etiquetas clínicas.

% Generado por Code/tools/tablas_anexos.py desde X5_auditoria.csv. No editar a mano.
\begin{table}[H]
\centering
\small
\caption{Auditoría metodológica independiente de consistencia, datos y modelos (Experimento X5).}
\label{tab:anexo_x5_auditoria}
\begin{adjustbox}{max width=\linewidth}
\begin{tabular}{p{11.5cm} r c}
\toprule
Propiedad auditada & Discrepancia máxima & Veredicto \\
\midrule
\multicolumn{3}{l}{\textbf{Integridad de etiquetas y particiones}} \\
\quad mass: every lesion of Module 4 exists in the original CSV files & 0.0 & Superada \\
\quad mass: labels equal MALIGNANT = 1 from the CSV files & 0.0 & Superada \\
\quad mass: train/test split equals the file of origin & 0.0 & Superada \\
\quad calcification: every lesion of Module 4 exists in the original CSV files & 0.0 & Superada \\
\quad calcification: labels equal MALIGNANT = 1 from the CSV files & 0.0 & Superada \\
\quad calcification: train/test split equals the file of origin & 0.0 & Superada \\
\quad mass: class counts equal the EDA (912 benign, 784 malignant) & 0.0 & Superada \\
\quad calcification: class counts equal the EDA (1199 benign, 673 malignant) & 0.0 & Superada \\
\multicolumn{3}{l}{\textbf{Ausencia de fuga de datos (escalado y folds)}} \\
\quad mass: angles refitted on train only equal Module 4, test included & 0.0 & Superada \\
\quad mass: every angle lies in [0, $\pi$], test included & 0.0 & Superada \\
\quad mass: no patient split across folds nor across train and test & 0.0 & Superada \\
\quad calcification: angles refitted on train only equal Module 4, test included & 0.0 & Superada \\
\quad calcification: every angle lies in [0, $\pi$], test included & 0.0 & Superada \\
\quad calcification: no patient split across folds nor across train and test & 0.0 & Superada \\
\multicolumn{3}{l}{\textbf{Equivalencia física de kernels cuánticos}} \\
\quad mass: C4 compute-uncompute kernel equals statevector kernel (12 x 12) & 0.0 & Superada \\
\quad mass: C5 compute-uncompute kernel equals statevector kernel (12 x 12) & 0.0 & Superada \\
\quad calcification: C4 compute-uncompute kernel equals statevector kernel (12 x 12) & 0.0 & Superada \\
\quad calcification: C5 compute-uncompute kernel equals statevector kernel (12 x 12) & 0.0 & Superada \\
\multicolumn{3}{l}{\textbf{Equivalencia de formulación QSVM (Qiskit vs dual)}} \\
\quad mass: QSVC of qiskit-machine-learning equals the precomputed-kernel SVM of X4 & 0.0 & Superada \\
\quad calcification: QSVC of qiskit-machine-learning equals the precomputed-kernel SVM of X4 & 0.0 & Superada \\
\multicolumn{3}{l}{\textbf{Invarianza de rotación y transformaciones}} \\
\quad changing x\_i changes only <Z\_i> of C5, to -sin(2 x\_i) & 0.0 & Superada \\
\quad C4: the kernel is invariant under a reordering of the features & 0.0 & Superada \\
\quad C5: the kernel is invariant under a reordering of the features & 0.0 & Superada \\
\multicolumn{3}{l}{\textbf{Consistencia de modelos de referencia clásicos}} \\
\quad mass: three classical models agree within one fold standard deviation & 2.39e-02 & Superada \\
\quad calcification: three classical models agree within one fold standard deviation & 1.63e-02 & Superada \\
\bottomrule
\end{tabular}
\end{adjustbox}
\par\smallskip
\parbox{0.95\linewidth}{\footnotesize Auditoría metodológica independiente de extremo a extremo (Experimento X5). Comprueba rigurosamente la ausencia de fugas de datos entre particiones de pacientes, la correspondencia exacta de etiquetas patológicas con los metadatos clínicos de CBIS-DDSM y la consistencia matemática de los componentes cuánticos.}
\end{table}

